\subsection{ULTRAFILTERS}

DEFINITION 4. An ultrafilter on a set X is a filter F such that there is no filter on X which is strictly finer than F (in other words, a maximal element in the ordered set of all filters on X).

Since the ordered set of all filters on X is inductive (no. 2, Proposition 1, Corollary 3), Zorn's lemma (Set Theory, R, § 6, no. 10) shows that:

THEOREM 1. If F is any filter on a set X, there is an ultrafilter finer than F.

PROPOSITION 5. Let \(\mathfrak{F}\) be an ultrafilter on a set \(\mathbf{X}\) . If \(\mathbf{A}\) and \(\mathbf{B}\) are two subsets of \(\mathbf{X}\) such that \(\mathbf{A} \cup \mathbf{B} \in \mathfrak{F}\) , then either \(\mathbf{A} \in \mathfrak{F}\) or \(\mathbf{B} \in \mathfrak{F}\) .

If the proposition is false, there exist subsets A and B of X such that A∉F and B∉F and A∪B∈F. Let G be the set of subsets M of X such that A∪M∈F. It is straightforward to check that G is a filter on X, and G is strictly finer than F, since B∈G; but this contradicts the hypothesis that F is an ultrafilter.

COROLLARY. If the union of a finite sequence \((\mathbf{A}_i)_{1 \leqslant i \leqslant n}\) of subsets of \(\mathbf{X}\) belongs to an ultrafilter \(\mathfrak{F}\) , then at least one of the \(\mathbf{A}_i\) belongs to \(\mathfrak{F}\) .

Proof is by induction on n.

In particular, if \((\mathbf{A}_i)_{1\leqslant i\leqslant n}\) is a covering of \(\mathbf{X}\) , then at least one of the \(\mathbf{A}_i\) belongs to \(\mathfrak{F}\) .

Proposition 5 characterizes the ultrafilters; more generally, we have:

PROPOSITION 6. Let \(\mathfrak{G}\) be a subbase of a filter on a set X. If for each subset Y of X we have either \(Y \in \mathfrak{G}\) or \(\complement Y \in \mathfrak{G}\) , then \(\mathfrak{G}\) is an ultrafilter on X.

Let \(\mathfrak{F}\) be a filter containing \(\mathfrak{G}\) (there is one, by hypothesis); then \(\mathfrak{F}\) coincides with \(\mathfrak{G}\) ; for if \(Y \in \mathfrak{F}\) then \(\complement Y \notin \mathfrak{F}\) ; hence \(\complement Y \notin \mathfrak{G}\) and therefore \(Y \in \mathfrak{G}\) .

Example of an ultrafilter. The set of all subsets of a non-empty set X which contain a given element \(a \in X\) is an ultrafilter; for it is a filter, and if Y is any subset of X then either \(a \in Y\) or \(a \in \complement Y\) . Such ultrafilters are called trivial.

Apart from this example, we shall never prove the existence of an ultrafilter (even on a countably infinite set) except by using Theorem 1 (and therefore the axiom of choice).

Remark. If X contains at least two elements, there are at least two distinct ultrafilters on X, and therefore the ordered set of filters on X has no greatest element.

PROPOSITION 7. Every filter \(\mathfrak{F}\) on a set \(X\) is the intersection of the ultrafilters finer than \(\mathfrak{F}\) .

Clearly this intersection contains \(\mathfrak{F}\) . Conversely, let \(A\) be a subset of \(X\) which does not belong to \(\mathfrak{F}\) , and let \(A'\) denote \(\complement A\) ; \(A\) contains no set of \(\mathfrak{F}\) ; hence every \(M \in \mathfrak{F}\) meets \(A'\) and therefore (no. 2, Proposition 1, Corollary 1) there is a filter \(\mathfrak{F}'\) which is finer than \(\mathfrak{F}\) and contains \(A'\) . If \(\mathfrak{U}\) is an ultrafilter finer than \(\mathfrak{F}'\) (Theorem 1) it follows that \(A \notin \mathfrak{U}\) . This completes the proof.

